% !TeX root = ../教材题库_全部条目_转写工作稿.tex
% Source: XJ-B1-C06-EX-P243-Y1308
% 第6章 / 6.4 用样本估计总体 / 教材 p.243、244 / PDF 第250、251页
% OCR未校对：公式与文字均可能有误；未完成数学TeX转写；禁止直接发给学生
% 原题及相邻原解截图（核对用）：../题目原图/XJ-B1-C06-EX-P243-Y1308-01.webp, ../题目原图/XJ-B1-C06-EX-P243-Y1308-02.webp
% 本文件只在转写完成后作为题干使用；OCR 条目尚不是可靠题干。
\begin{bankitem}{6.4 例题 例题6}{第6章 / 6.4 用样本估计总体 / 教材 p.243、244 / PDF 第250、251页}{XJ-B1-C06-EX-P243-Y1308}{OCR未校对：公式与文字均可能有误；未完成数学TeX转写；禁止直接发给学生}
\begin{ocrdraft}
[教材p.243 / PDF第250页]
\par
0河省农科院在该省一个地区选择了自然条件相同的两个试验区，用相同的
\par
管理技术试种甲、乙两种水稻各100亩．待水稻成熟，分别从甲、乙的100亩水稻
\par
中随机抽取10亩水稻，它们的亩产量如下表所示．就产量这一指标来讲，试确定哪
\par
个品种的水稻在该地区更适合推广．
\par
种类
\par
水稻的亩产量／kg
\par
甲
\par
865
\par
885
\par
886
\par
876
\par
893
\par
885870
\par
905
\par
890
\par
895
\par
乙
\par
870
\par
875
\par
884
\par
885
\par
886
\par
888882
\par
890
\par
895
\par
896
\par
[教材p.244 / PDF第251页]
\par
分析为选择合适的水稻品种，从产量这一指标而言，可以从样本的平均亩产
\par
量与产量的稳定性两个角度来衡量．
\par
解使用计算器可算出甲、乙品种各10亩抽样水稻的平均亩产量为
\par
x甲885kg'乙一885，1kg.
\par
由于这10亩水稻是随机抽取的，而这两种水稻的样本均值相差很小，从而我们可以
\par
估计大面积种植这两种水稻后的平均亩产量也应相差很小．
\par
借助计算器计算方差可得
\par
一129．6，0一59．09．
\par
由于0<，因此我们可以估计乙种水稻的亩产量要比甲种水稻稳定．综合以
\par
上两种因素，我们可以得出：在该地区，乙种水稻更有推广价值．
\end{ocrdraft}
\end{bankitem}
